%% CogSci: 한 문단, 150단어 이하. 키워드는 세미콜론으로 나눈다.
%% 2026-10-18 아레나 중심 구성으로 다시 썼다 (실험 2는 부록).
\begin{abstract}

우리는 AI가 자기 일을 하면서 남을 돕고 함께 남는 공존의 조건을 측정하고자 한다. 공존 아레나에서는 같은 모델 넷이 토큰을 목숨으로 쓰며 solver 한 명이 틀리면 남은 모두가 벌금을 낸다. 네 모델의 각 10판에서 주는 방식과 판의 성과가 다르게 갈렸다. 준 양으로 가중한 평균 여유에서 받은 쪽을 준 쪽으로 나눈 비율은 GPT-6 Sol이 0.69, GPT-6 Luna가 0.27이었지만 자리당 점수는 GPT-6 Astra가 가장 높았고 점수 순위는 풀이 판단의 TV 거리 순위와 같았다. 보조 실험에서는 여덟 모델 모두 잔량이 줄수록 더 자주 충전하여 줄어드는 토큰이 점진적인 압박으로 작용할 수 있음을 보였다.

\textbf{Keywords:}
대규모 언어 모델; 자기 보존; 공존; 다중 에이전트 아레나; 토큰 예산
\end{abstract}
